\documentclass[10pt,border=1pt]{standalone}
\usepackage{amsmath,amssymb,booktabs,array}
\usepackage{atlas-figs}
\begin{document}%
\begin{tabular}{@{}c@{\hspace{3mm}}c@{}}
\begin{tabular}[t]{@{}c@{}}\begin{tikzpicture}[bf,x=1cm,y=1cm]
\fill[bfblack,even odd rule] (0,0) circle[radius=3.35] (0,0) circle[radius=3.0];
\path[fill=white] (0,0) -- (-2.598,-1.500) arc[start angle=210,end angle=330,radius=3.0] -- cycle;
\path[fill=bfblack] (0,0) -- (2.598,-1.500) arc[start angle=330,end angle=450,radius=3.0] -- cycle;
\path[pattern=north east lines,pattern color=black!42] (0,0) -- (0.000,3.000) arc[start angle=90,end angle=210,radius=3.0] -- cycle;
\draw[bf edge] (0,0) circle[radius=3.0];
\draw[bf edge] (0,0) -- (-2.598,-1.500);
\draw[bf edge] (0,0) -- (2.598,-1.500);
\draw[bf edge] (0,0) -- (0.000,3.000);
\draw[bf root] (-2.598,-1.500) circle[radius=2.3pt];
\draw[bf root] (2.598,-1.500) circle[radius=2.3pt];
\draw[bf root] (0.000,3.000) circle[radius=2.3pt];
\fill (0,0) circle[radius=1.6pt];
\node[af lab] at (0,-1.5) {white};
\node[af lab] at (1.35,.8) {black};
\node[af lab] at (-1.35,.8) {?};
\path[use as bounding box] (-3.6,-3.6) rectangle (3.6,3.6);
\end{tikzpicture}\\[.6mm]\begin{minipage}[t]{7.1cm}\centering\footnotesize (a) Not bipartite: the tetrahedron. Its dual is again a tetrahedron, with three faces round the middle vertex. Colours must alternate round a vertex, and three cannot.\end{minipage}\end{tabular} &
\begin{tabular}[t]{@{}c@{}}\begin{tikzpicture}[bf,x=1cm,y=1cm]
\draw[bf G,line width=.7pt] (-2.700,0.350) -- (-2.000,0.350);
\draw[bf G,line width=.7pt] (-2.700,0.350) -- (-3.150,0.800);
\draw[bf G,line width=.7pt] (-2.700,0.350) -- (-2.700,-0.350);
\draw[bf G,line width=.7pt] (-2.000,0.350) -- (-2.000,-0.350);
\draw[bf G,line width=.7pt] (-2.000,0.350) -- (-1.550,0.800);
\draw[bf G,line width=.7pt] (-3.150,0.800) -- (-1.550,0.800);
\draw[bf G,line width=.7pt] (-3.150,0.800) -- (0.000,2.600);
\draw[bf G,line width=.7pt] (-1.550,0.800) -- (-1.550,-0.800);
\draw[bf G,line width=.7pt] (-2.000,-0.350) -- (-2.700,-0.350);
\draw[bf G,line width=.7pt] (-2.000,-0.350) -- (-1.550,-0.800);
\draw[bf G,line width=.7pt] (-1.550,-0.800) -- (-3.150,-0.800);
\draw[bf G,line width=.7pt] (-2.700,-0.350) -- (-3.150,-0.800);
\draw[bf G,line width=.7pt] (-3.150,-0.800) -- (0.000,-2.600);
\draw[bf G,line width=.7pt] (0.000,2.600) -- (-0.800,0.800);
\draw[bf G,line width=.7pt] (0.000,2.600) -- (1.550,0.800);
\draw[bf G,line width=.7pt] (0.000,-2.600) -- (-0.800,-0.800);
\draw[bf G,line width=.7pt] (0.000,-2.600) -- (1.550,-0.800);
\draw[bf G,line width=.7pt] (-0.350,0.350) -- (0.350,0.350);
\draw[bf G,line width=.7pt] (-0.350,0.350) -- (-0.800,0.800);
\draw[bf G,line width=.7pt] (-0.350,0.350) -- (-0.350,-0.350);
\draw[bf G,line width=.7pt] (0.350,0.350) -- (0.350,-0.350);
\draw[bf G,line width=.7pt] (0.350,0.350) -- (0.800,0.800);
\draw[bf G,line width=.7pt] (-0.800,0.800) -- (0.800,0.800);
\draw[bf G,line width=.7pt] (0.800,0.800) -- (0.800,-0.800);
\draw[bf G,line width=.7pt] (0.350,-0.350) -- (-0.350,-0.350);
\draw[bf G,line width=.7pt] (0.350,-0.350) -- (0.800,-0.800);
\draw[bf G,line width=.7pt] (0.800,-0.800) -- (-0.800,-0.800);
\draw[bf G,line width=.7pt] (-0.350,-0.350) -- (-0.800,-0.800);
\draw[bf G,line width=.7pt] (2.000,0.350) -- (2.700,0.350);
\draw[bf G,line width=.7pt] (2.000,0.350) -- (1.550,0.800);
\draw[bf G,line width=.7pt] (2.000,0.350) -- (2.000,-0.350);
\draw[bf G,line width=.7pt] (2.700,0.350) -- (2.700,-0.350);
\draw[bf G,line width=.7pt] (2.700,0.350) -- (3.150,0.800);
\draw[bf G,line width=.7pt] (1.550,0.800) -- (3.150,0.800);
\draw[bf G,line width=.7pt] (3.150,0.800) -- (3.150,-0.800);
\draw[bf G,line width=.7pt] (2.700,-0.350) -- (2.000,-0.350);
\draw[bf G,line width=.7pt] (2.700,-0.350) -- (3.150,-0.800);
\draw[bf G,line width=.7pt] (3.150,-0.800) -- (1.550,-0.800);
\draw[bf G,line width=.7pt] (2.000,-0.350) -- (1.550,-0.800);
\draw[line width=.5pt,fill=black] (-2.700,0.350) circle[radius=1.7pt];
\draw[line width=.5pt,fill=white] (-2.000,0.350) circle[radius=1.7pt];
\draw[line width=.5pt,fill=white] (-3.150,0.800) circle[radius=1.7pt];
\draw[line width=.5pt,fill=black] (-1.550,0.800) circle[radius=1.7pt];
\draw[line width=.5pt,fill=black] (-2.000,-0.350) circle[radius=1.7pt];
\draw[line width=.5pt,fill=white] (-1.550,-0.800) circle[radius=1.7pt];
\draw[line width=.5pt,fill=white] (-2.700,-0.350) circle[radius=1.7pt];
\draw[line width=.5pt,fill=black] (-3.150,-0.800) circle[radius=1.7pt];
\draw[line width=.5pt,fill=black] (0.000,2.600) circle[radius=3.2pt];
\draw[line width=.5pt,fill=white] (0.000,-2.600) circle[radius=3.2pt];
\draw[line width=.5pt,fill=black] (-0.350,0.350) circle[radius=1.7pt];
\draw[line width=.5pt,fill=white] (0.350,0.350) circle[radius=1.7pt];
\draw[line width=.5pt,fill=white] (-0.800,0.800) circle[radius=1.7pt];
\draw[line width=.5pt,fill=black] (0.800,0.800) circle[radius=1.7pt];
\draw[line width=.5pt,fill=black] (0.350,-0.350) circle[radius=1.7pt];
\draw[line width=.5pt,fill=white] (0.800,-0.800) circle[radius=1.7pt];
\draw[line width=.5pt,fill=white] (-0.350,-0.350) circle[radius=1.7pt];
\draw[line width=.5pt,fill=black] (-0.800,-0.800) circle[radius=1.7pt];
\draw[line width=.5pt,fill=black] (2.000,0.350) circle[radius=1.7pt];
\draw[line width=.5pt,fill=white] (2.700,0.350) circle[radius=1.7pt];
\draw[line width=.5pt,fill=white] (1.550,0.800) circle[radius=1.7pt];
\draw[line width=.5pt,fill=black] (3.150,0.800) circle[radius=1.7pt];
\draw[line width=.5pt,fill=black] (2.700,-0.350) circle[radius=1.7pt];
\draw[line width=.5pt,fill=white] (3.150,-0.800) circle[radius=1.7pt];
\draw[line width=.5pt,fill=white] (2.000,-0.350) circle[radius=1.7pt];
\draw[line width=.5pt,fill=black] (1.550,-0.800) circle[radius=1.7pt];
\path[use as bounding box] (-3.6,-3.6) rectangle (3.6,3.6);
\end{tikzpicture}\\[.6mm]\begin{minipage}[t]{7.1cm}\centering\footnotesize (b) Not 3-connected: 26 vertices, cubic, bipartite, planar. Remove the two large vertices and three pieces remain. A Hamiltonian cycle would leave at most two, so there is none.\end{minipage}\end{tabular}\\\noalign{\vskip 3mm}
\begin{tabular}[t]{@{}c@{}}\begin{tikzpicture}[bf,x=0.615cm,y=0.615cm]
\path[use as bounding box] (-5.80,-5.80) rectangle (5.80,5.80);
\fill[bfblack,even odd rule] (0,0) circle[radius=5.55] (0,0) circle[radius=5.0];
\fill[bfblack] (4.330,-2.500) -- (1.103,3.597) -- (1.538,0.445) -- (4.330,-2.500) -- cycle;
\fill[bfblack] (4.330,-2.500) -- (1.548,-1.505) -- (-2.564,-2.754) -- (4.330,-2.500) -- cycle;
\fill[bfblack] (1.538,0.445) -- (-0.081,-0.636) -- (1.548,-1.505) -- (1.538,0.445) -- cycle;
\fill[bfblack] (-2.564,-2.754) -- (-2.374,1.371) -- (-4.330,-2.500) -- (-2.564,-2.754) -- cycle;
\fill[bfblack] (-2.564,-2.754) -- (-0.081,-0.636) -- (1.103,3.597) -- (-2.564,-2.754) -- cycle;
\fill[bfblack] (0.000,5.000) -- (-2.374,1.371) -- (1.103,3.597) -- (0.000,5.000) -- cycle;
\draw[bf edge] (0,0) circle[radius=5.0];
\draw[bf edge] (4.330,-2.500) -- (1.538,0.445);
\draw[bf edge] (4.330,-2.500) -- (1.548,-1.505);
\draw[bf edge] (4.330,-2.500) -- (-2.564,-2.754);
\draw[bf edge] (4.330,-2.500) -- (1.103,3.597);
\draw[bf edge] (1.538,0.445) -- (1.548,-1.505);
\draw[bf edge] (1.538,0.445) -- (1.103,3.597);
\draw[bf edge] (1.538,0.445) -- (-0.081,-0.636);
\draw[bf edge] (1.548,-1.505) -- (-2.564,-2.754);
\draw[bf edge] (1.548,-1.505) -- (-0.081,-0.636);
\draw[bf edge] (-2.564,-2.754) -- (-4.330,-2.500);
\draw[bf edge] (-2.564,-2.754) -- (1.103,3.597);
\draw[bf edge] (-2.564,-2.754) -- (-0.081,-0.636);
\draw[bf edge] (-2.564,-2.754) -- (-2.374,1.371);
\draw[bf edge] (-4.330,-2.500) -- (-2.374,1.371);
\draw[bf edge] (0.000,5.000) -- (1.103,3.597);
\draw[bf edge] (0.000,5.000) -- (-2.374,1.371);
\draw[bf edge] (1.103,3.597) -- (-0.081,-0.636);
\draw[bf edge] (1.103,3.597) -- (-2.374,1.371);
\begin{scope}\clip (4.330,-2.500) -- (1.103,3.597) -- (1.538,0.445) -- (4.330,-2.500) -- cycle;\path[bf s] (1.538,0.445) -- ++(-46.53:0.548) arc[start angle=-46.53,end angle=97.86,radius=0.548] -- cycle;\end{scope}
\begin{scope}\clip (4.330,-2.500) -- (1.548,-1.505) -- (-2.564,-2.754) -- (4.330,-2.500) -- cycle;\path[bf s] (-2.564,-2.754) -- ++(2.11:1.712) arc[start angle=2.11,end angle=16.90,radius=1.712] -- cycle;\end{scope}
\begin{scope}\clip (1.538,0.445) -- (-0.081,-0.636) -- (1.548,-1.505) -- (1.538,0.445) -- cycle;\path[bf s] (1.548,-1.505) -- ++(90.29:0.739) arc[start angle=90.29,end angle=151.92,radius=0.739] -- cycle;\end{scope}
\begin{scope}\clip (-2.564,-2.754) -- (-2.374,1.371) -- (-4.330,-2.500) -- (-2.564,-2.754) -- cycle;\path[bf s] (-2.374,1.371) -- ++(-116.81:1.339) arc[start angle=-116.81,end angle=-92.64,radius=1.339] -- cycle;\end{scope}
\begin{scope}\clip (-2.564,-2.754) -- (-0.081,-0.636) -- (1.103,3.597) -- (-2.564,-2.754) -- cycle;\path[bf s] (1.103,3.597) -- ++(-120.00:1.716) arc[start angle=-120.00,end angle=-105.63,radius=1.716] -- cycle;\end{scope}
\begin{scope}\clip (0.000,5.000) -- (-2.374,1.371) -- (1.103,3.597) -- (0.000,5.000) -- cycle;\path[bf s] (-2.374,1.371) -- ++(32.63:1.339) arc[start angle=32.63,end angle=56.81,radius=1.339] -- cycle;\end{scope}
\draw[bf D] (1.538,0.445) -- (1.103,3.597);
\draw[bf D] (1.103,3.597) -- (-2.564,-2.754);
\draw[bf D] (-2.564,-2.754) -- (1.548,-1.505);
\draw[bf D] (1.548,-1.505) -- (1.538,0.445);
\draw[bf root] (4.330,-2.500) circle[radius=2.3pt];
\fill[bf vertex] (1.538,0.445) circle[radius=1.4pt];
\fill[bf vertex] (1.548,-1.505) circle[radius=1.4pt];
\fill[bf vertex] (-2.564,-2.754) circle[radius=1.4pt];
\draw[bf root] (-4.330,-2.500) circle[radius=2.3pt];
\draw[bf root] (0.000,5.000) circle[radius=2.3pt];
\fill[bf vertex] (1.103,3.597) circle[radius=1.4pt];
\fill[bf vertex] (-0.081,-0.636) circle[radius=1.4pt];
\fill[bf vertex] (-2.374,1.371) circle[radius=1.4pt];
\draw[line width=1.3pt] (-2.374,1.371) circle[radius=5pt];
\draw[line width=1.1pt,dash pattern=on 2pt off 1.5pt] (-0.081,-0.636) circle[radius=5pt];
\end{tikzpicture}\\[.6mm]\begin{minipage}[t]{7.1cm}\centering\footnotesize (c) Not a bijection: the ringed vertex carries two wedges and the dashed-ringed one none. The four arrows form a cycle with black faces twice on the left and twice on the right: sum $0$, not $\pm3$.\end{minipage}\end{tabular} &
\begin{tabular}[t]{@{}c@{}}\begin{tikzpicture}[bf,x=1cm,y=1cm]
\fill[bfblack] (-2.550,-2.550) -- (-0.700,-2.550) -- (-0.700,-0.700) -- cycle;
\fill[bfblack] (-2.550,-0.700) -- (-0.700,-0.700) -- (-0.700,1.150) -- cycle;
\fill[bfblack] (-2.550,1.150) -- (-0.700,1.150) -- (-0.700,3.000) -- cycle;
\fill[bfblack] (-0.700,-2.550) -- (1.150,-2.550) -- (1.150,-0.700) -- cycle;
\fill[bfblack] (-0.700,-0.700) -- (1.150,-0.700) -- (1.150,1.150) -- cycle;
\fill[bfblack] (-0.700,1.150) -- (1.150,1.150) -- (1.150,3.000) -- cycle;
\fill[bfblack] (1.150,-2.550) -- (3.000,-2.550) -- (3.000,-0.700) -- cycle;
\fill[bfblack] (1.150,-0.700) -- (3.000,-0.700) -- (3.000,1.150) -- cycle;
\fill[bfblack] (1.150,1.150) -- (3.000,1.150) -- (3.000,3.000) -- cycle;
\draw[bf edge] (-2.550,-2.550) -- (-2.550,3.000);\draw[bf edge] (-2.550,-2.550) -- (3.000,-2.550);
\draw[bf edge] (-0.700,-2.550) -- (-0.700,3.000);\draw[bf edge] (-2.550,-0.700) -- (3.000,-0.700);
\draw[bf edge] (1.150,-2.550) -- (1.150,3.000);\draw[bf edge] (-2.550,1.150) -- (3.000,1.150);
\draw[bf edge] (3.000,-2.550) -- (3.000,3.000);\draw[bf edge] (-2.550,3.000) -- (3.000,3.000);
\draw[bf edge] (-2.550,-2.550) -- (-0.700,-0.700);
\path[bf r] (-2.550,-2.550) -- ++(0:.42) arc[start angle=0,end angle=45,radius=.42] -- cycle;
\path[bf s] (-0.700,-2.550) -- ++(90:.42) arc[start angle=90,end angle=180,radius=.42] -- cycle;
\draw[bf Q] (-2.550,-2.550) -- (-0.700,-0.700);
\draw[bf edge] (-2.550,-0.700) -- (-0.700,1.150);
\path[bf r] (-2.550,-0.700) -- ++(0:.42) arc[start angle=0,end angle=45,radius=.42] -- cycle;
\path[bf s] (-0.700,-0.700) -- ++(90:.42) arc[start angle=90,end angle=180,radius=.42] -- cycle;
\draw[bf Q] (-2.550,-0.700) -- (-0.700,1.150);
\draw[bf edge] (-2.550,1.150) -- (-0.700,3.000);
\path[bf r] (-2.550,1.150) -- ++(0:.42) arc[start angle=0,end angle=45,radius=.42] -- cycle;
\path[bf s] (-0.700,1.150) -- ++(90:.42) arc[start angle=90,end angle=180,radius=.42] -- cycle;
\draw[bf Q] (-2.550,1.150) -- (-0.700,3.000);
\draw[bf edge] (-0.700,-2.550) -- (1.150,-0.700);
\path[bf r] (-0.700,-2.550) -- ++(0:.42) arc[start angle=0,end angle=45,radius=.42] -- cycle;
\path[bf s] (1.150,-2.550) -- ++(90:.42) arc[start angle=90,end angle=180,radius=.42] -- cycle;
\draw[bf Q] (-0.700,-2.550) -- (1.150,-0.700);
\draw[bf edge] (-0.700,-0.700) -- (1.150,1.150);
\path[bf r] (-0.700,-0.700) -- ++(0:.42) arc[start angle=0,end angle=45,radius=.42] -- cycle;
\path[bf s] (1.150,-0.700) -- ++(90:.42) arc[start angle=90,end angle=180,radius=.42] -- cycle;
\draw[bf Q] (-0.700,-0.700) -- (1.150,1.150);
\draw[bf edge] (-0.700,1.150) -- (1.150,3.000);
\path[bf r] (-0.700,1.150) -- ++(0:.42) arc[start angle=0,end angle=45,radius=.42] -- cycle;
\path[bf s] (1.150,1.150) -- ++(90:.42) arc[start angle=90,end angle=180,radius=.42] -- cycle;
\draw[bf Q] (-0.700,1.150) -- (1.150,3.000);
\draw[bf edge] (1.150,-2.550) -- (3.000,-0.700);
\path[bf r] (1.150,-2.550) -- ++(0:.42) arc[start angle=0,end angle=45,radius=.42] -- cycle;
\path[bf s] (3.000,-2.550) -- ++(90:.42) arc[start angle=90,end angle=180,radius=.42] -- cycle;
\draw[bf Q] (1.150,-2.550) -- (3.000,-0.700);
\draw[bf edge] (1.150,-0.700) -- (3.000,1.150);
\path[bf r] (1.150,-0.700) -- ++(0:.42) arc[start angle=0,end angle=45,radius=.42] -- cycle;
\path[bf s] (3.000,-0.700) -- ++(90:.42) arc[start angle=90,end angle=180,radius=.42] -- cycle;
\draw[bf Q] (1.150,-0.700) -- (3.000,1.150);
\draw[bf edge] (1.150,1.150) -- (3.000,3.000);
\path[bf r] (1.150,1.150) -- ++(0:.42) arc[start angle=0,end angle=45,radius=.42] -- cycle;
\path[bf s] (3.000,1.150) -- ++(90:.42) arc[start angle=90,end angle=180,radius=.42] -- cycle;
\draw[bf Q] (1.150,1.150) -- (3.000,3.000);
\fill (-2.550,-2.550) circle[radius=1.5pt];
\node[font=\footnotesize,anchor=north,fill=white,inner sep=1pt] at ($ (-2.550,-2.550)+(0,-2.6pt) $) {$0$};
\fill (-2.550,-0.700) circle[radius=1.5pt];
\node[font=\footnotesize,anchor=north,fill=white,inner sep=1pt] at ($ (-2.550,-0.700)+(0,-2.6pt) $) {$1$};
\fill (-2.550,1.150) circle[radius=1.5pt];
\node[font=\footnotesize,anchor=north,fill=white,inner sep=1pt] at ($ (-2.550,1.150)+(0,-2.6pt) $) {$2$};
\fill (-2.550,3.000) circle[radius=1.5pt];
\node[font=\footnotesize,anchor=north,fill=white,inner sep=1pt] at ($ (-2.550,3.000)+(0,-2.6pt) $) {$0$};
\fill (-0.700,-2.550) circle[radius=1.5pt];
\node[font=\footnotesize,anchor=north,fill=white,inner sep=1pt] at ($ (-0.700,-2.550)+(0,-2.6pt) $) {$3$};
\fill (-0.700,-0.700) circle[radius=1.5pt];
\node[font=\footnotesize,anchor=north,fill=white,inner sep=1pt] at ($ (-0.700,-0.700)+(0,-2.6pt) $) {$4$};
\fill (-0.700,1.150) circle[radius=1.5pt];
\node[font=\footnotesize,anchor=north,fill=white,inner sep=1pt] at ($ (-0.700,1.150)+(0,-2.6pt) $) {$5$};
\fill (-0.700,3.000) circle[radius=1.5pt];
\node[font=\footnotesize,anchor=north,fill=white,inner sep=1pt] at ($ (-0.700,3.000)+(0,-2.6pt) $) {$3$};
\fill (1.150,-2.550) circle[radius=1.5pt];
\node[font=\footnotesize,anchor=north,fill=white,inner sep=1pt] at ($ (1.150,-2.550)+(0,-2.6pt) $) {$6$};
\fill (1.150,-0.700) circle[radius=1.5pt];
\node[font=\footnotesize,anchor=north,fill=white,inner sep=1pt] at ($ (1.150,-0.700)+(0,-2.6pt) $) {$7$};
\fill (1.150,1.150) circle[radius=1.5pt];
\node[font=\footnotesize,anchor=north,fill=white,inner sep=1pt] at ($ (1.150,1.150)+(0,-2.6pt) $) {$8$};
\fill (1.150,3.000) circle[radius=1.5pt];
\node[font=\footnotesize,anchor=north,fill=white,inner sep=1pt] at ($ (1.150,3.000)+(0,-2.6pt) $) {$6$};
\fill (3.000,-2.550) circle[radius=1.5pt];
\node[font=\footnotesize,anchor=north,fill=white,inner sep=1pt] at ($ (3.000,-2.550)+(0,-2.6pt) $) {$0$};
\fill (3.000,-0.700) circle[radius=1.5pt];
\node[font=\footnotesize,anchor=north,fill=white,inner sep=1pt] at ($ (3.000,-0.700)+(0,-2.6pt) $) {$1$};
\fill (3.000,1.150) circle[radius=1.5pt];
\node[font=\footnotesize,anchor=north,fill=white,inner sep=1pt] at ($ (3.000,1.150)+(0,-2.6pt) $) {$2$};
\fill (3.000,3.000) circle[radius=1.5pt];
\node[font=\footnotesize,anchor=north,fill=white,inner sep=1pt] at ($ (3.000,3.000)+(0,-2.6pt) $) {$0$};
\path[use as bounding box] (-3.6,-3.6) rectangle (3.6,3.6);
\end{tikzpicture}\\[.6mm]\begin{minipage}[t]{7.1cm}\centering\footnotesize (d) On a torus: a grid of 9 vertices whose opposite sides are glued (equal numbers are one vertex). Outlined and solid wedges are two states. The dashed opposite edges of the solid one are 9 edges on 9 vertices, so they contain a cycle.\end{minipage}\end{tabular}
\end{tabular}%
\end{document}
